% Physical pages 3 and 4 of the concise study: the thematic orientation, in the
% order of the Mass, with what each text says, keeps and omits. Every claim and
% quotation comes from the reviewed expansive study. The developed comparison
% of the three interpretations begins on the fifth page.

\runninghead{Themes and Movement}
\section{The Propers: Themes and Movement}
\zlabel{triptych:concise:themes:start}

The Gospel of this last Sunday after Pentecost is the second half of the
answer Jesus gave on the Mount of Olivet, when the disciples asked when the
temple would be destroyed and what would be the sign of his coming ``and of
the consummation of the world'' (Mt 24:3). Three interpretations of the Mass
take up three questions that its words raise. The first asks what the end is
a gathering to. The Gospel says, ``Wheresoever the body shall be, there shall
the eagles also be gathered together'' (24:28), and St Jerome, St Hilary of
Poitiers and St Augustine each take the body for Christ's; this
interpretation finds the gathering already begun in the Introit's return of
the captivity ``from all places'' and in the Epistle's thanks to the Father
who has ``translated us into the kingdom of the Son of his love.'' The second
asks to whom the commands to flee from Judea are spoken. St John Chrysostom
answers that they concern the Roman war and the Jews of that time. Hilary
reads the command to the man on the housetop of the person reborn, who must
not go down again into the desire of worldly things, and Jerome adds to the
plain senses a flight from every perverse doctrine; this interpretation takes
the Secret's petition to be ``delivered from earthly desires'' for the prayer
of those who keep that rule. The third asks what makes a promise certain when
no one knows its day. Chrysostom and St Gregory the Great argue from
prophecies already fulfilled, and Augustine from God, with whom the future is
as certain as the past; this interpretation reads the past tenses of the
chants and the Communion's ``believe that you shall receive'' by their
answers.

The Introit's words come from the letter that Jeremias sent from Jerusalem to
the exiles carried away to Babylon. Before its promise the letter tells them
to build houses, plant orchards and ``seek the peace of the city, to which I
have caused you to be carried away captives'' (Jer 29:5, 7); it warns them
against the prophets and diviners among them, and puts God's visit after
seventy years (vv.~8--10). The antiphon keeps three things that God says he
will do: think thoughts of peace, hear those who call on him, and bring back
their captivity. It leaves out the promise ``to give you an end and
patience,'' the seeking ``with all your heart'' (v.~13) and ``I will gather
you out of all nations,'' so that ``from all places'' follows straight on ``I
will bring back your captivity.'' Its opening, \latin{Dicit Dóminus}, ``The
Lord saith,'' makes the whole of it God's speech in the first person. The
psalm verse speaks in the past tense, ``Thou hast turned away the captivity of
Jacob'' (Ps 84:2), though two verses later the same psalm still prays,
``Convert us, O God our saviour'' (v.~5).

The Collect asks God to ``arouse \dots\ the wills of Thy faithful,'' that,
``by more earnestly following after the fruit of the divine work, they may the
more abundantly partake of Thy mercies.'' Its Latin has two comparatives: the
faithful are to follow the fruit \latin{propénsius}, ``more readily,'' and to
receive \latin{remédia maióra}, ``greater remedies,'' which Lasance gives as
``mercies.'' \latin{Divíni óperis}, ``of the divine work,'' can be the work
that God does or the work done for God, and the prayer does not say which.

The Epistle is the prayer with which St Paul opens his letter to the
Colossians, a church taught by Epaphras, of whose faith and love he had heard
(Col 1:4, 7--8). The lesson is a single sentence. Paul asks that they be
filled with the knowledge of God's will and ``walk worthy of God, in all
things pleasing,'' bearing fruit, growing in the knowledge of God and
strengthened ``in all patience and longsuffering with joy''; and he thanks the
Father, who ``hath made us worthy to be partakers of the lot of the saints in
light,'' ``hath delivered us from the power of darkness'' and ``hath
translated us'' into his Son's kingdom, ``In whom we have redemption through
his blood'' (vv.~9--14). The petitions ask for what is not yet complete, and
the clauses of thanks name what has been done. On ``that ye may be filled''
Chrysostom notes that Paul says ``not, `that ye may receive,' for they had
received'' (\work{Hom.\ in Col.}~2).

The Gradual takes two verses from Psalm~43. The psalm remembers that the
fathers ``got not the possession of the land by their own sword'' (v.~4); it
then complains that God has cast his people off and given them up ``like
sheep to be eaten'' (vv.~10, 12), and near its end it cries, ``Arise, why
sleepest thou, O Lord?'' (v.~23). The chant sings the confidence and not the
complaint: \latin{Liberásti nos}, ``Thou hast delivered us,'' where the
Vulgate's verb is \latin{salvasti}, ``thou hast saved,'' and the promise to
glory in God all the day. Augustine, preaching on the whole psalm, takes its
complaint for the voice of persecuted Christians, who ask, ``Why did he will
us to suffer all these things in the meantime?'' (\work{Enarrationes in
Psalmos}~43.10). The Alleluia and the Offertory sing the opening of
Psalm~129, ``From the depths I have cried to Thee, O Lord: Lord, hear my
prayer,'' with the Missal's ``my prayer'' where the Vulgate has ``my voice.''
The Offertory repeats the first line, so that the cry is sung three times in
the Mass. The rest of the psalm, which is not sung, goes on to ``merciful
forgiveness'' and ``plentiful redemption'' (vv.~4, 7).

The Gospel begins at the fifteenth verse of the Lord's answer, and the Missal
drops the ``therefore'' that ties it to the disciples' question. It has four
parts: the sign from Daniel and the commands to flee, ending with the days
shortened ``for the sake of the elect'' (vv.~15--22); the false christs who
will be said to be ``in the desert'' or ``in the closets,'' the lightning,
and the body and the eagles (vv.~23--28); the darkened sky, the sign of the
Son of man, his coming on the clouds and the gathering of the elect ``from
the four winds'' (vv.~29--31); and the fig tree with two sayings, ``this
generation shall not pass till all these things be done'' and ``Heaven and
earth shall pass: but my words shall not pass'' (vv.~32--35). The next verse,
not read, says that no one knows the day and the hour (v.~36). Jerome takes
the parenthesis ``he that readeth let him understand'' as a summons to a
mystical sense, and gives three plain senses of the abomination without
choosing among them: Antichrist, the image of Caesar that Pilate placed in the
temple, or the equestrian statue of Hadrian (\work{In Matthaeum}~IV, on
24:15). Chrysostom takes it for ``the statue of him who then took the city''
(\work{Hom.\ in Mt.}~75.2). He says that the ``Then'' of verse~23 passes over
``all the intermediate time from the taking of Jerusalem unto the preludes of
the consummation'' (76.2), and Jerome that the false christs are ``better
understood of the consummation of the world'' (on 24:23). Both read verses
29--31 of the Lord's visible coming, and so do all three interpretations.

At the offering the Secret asks God to ``accept the offerings and prayers of
Thy people'' and to ``turn all our hearts unto Thee, that, being delivered
from earthly desires, we may pass on to the enjoyments of heaven.'' Lasance's
``enjoyments'' stands for \latin{desidéria}, so that the Latin names desire
on both sides: from earthly \latin{cupiditátibus} to heavenly
\latin{desidéria}. Its \latin{liberáti}, ``delivered,'' is the verb of the
Gradual's \latin{Liberásti}, and \latin{transeámus}, ``may we pass over,'' is
a verb of travel.

The Communion is the Lord's promise in Mark, given when Peter pointed out the
fig tree that Jesus had cursed and that had withered from the roots
overnight. Jesus answered, ``Have the faith of God,'' and said that whoever
``shall not stagger in his heart, but believe'' can bid a mountain be cast
into the sea (Mk 11:22--23). The antiphon sings the next verse, opening with
the ``Amen I say to you'' and closing with the ``it shall be done'' of the
verse before it; its promise is in the future tense, \latin{crédite quia
accipiétis}, ``believe that you shall receive.'' The verse after it adds,
``when you shall stand to pray, forgive, if you have aught against any man''
(11:25). The Postcommunion asks that, ``through this sacrament which we have
received, whatever is evil in our hearts may be restored by its gift of
healing.'' Its Latin speaks of the mind, \latin{mente}, and of a cure,
\latin{medicatiónis}, as the Collect asked for \latin{remédia}. Dom Lucien
Fromage, the Solesmes monk who continued Dom Guéranger's \work{The Liturgical
Year}, prays on this Sunday that the Sacrament may cure ``whatsoever there may
remain faulty in our souls, at this close of the Year'' (vol.~11, p.~513).
\zlabel{triptych:concise:themes:end}
